\section*{अध्याय \ref{ch:g9:thales} --- थेल्स प्रमेय}

\begin{solution}{exo:g9:thales:1}
थेल्स: $\dfrac{AM}{AB} = \dfrac{AN}{AC} = \dfrac{MN}{BC}$, यानी
$\dfrac26 = \dfrac13$। तब $\dfrac{3}{AC} = \dfrac13$ से $AC = 9$, और
$\dfrac{MN}{9} = \dfrac13$ से $MN = 3$।
\end{solution}

\begin{solution}{exo:g9:thales:2}
पहले $AB = AM + MB = 5 + 3 = 8$ निकालो। फिर
$\dfrac{AM}{AB} = \dfrac{AN}{AC}$ से $\dfrac58 = \dfrac{4}{AC}$, यानी
$5 \times AC = 32$ और $AC = 6.4$।
\end{solution}

\begin{solution}{exo:g9:thales:3}
तितली ठीक वैसे ही चलती है जैसे एक के भीतर एक वाला विन्यास:
$\dfrac{AM}{AB} = \dfrac{AN}{AC} = \dfrac{MN}{BC}$, यानी
$\dfrac{3}{7.5} = \dfrac25$। तब $\dfrac{2}{AC} = \dfrac25$ से
$AC = 5$, और $\dfrac{2.6}{BC} = \dfrac25$ से
$BC = \dfrac{2.6 \times 5}{2} = 6.5$।
\end{solution}

\begin{solution}{exo:g9:thales:4}
$\dfrac{AM}{AB} = \dfrac68 = \dfrac34$ और
$\dfrac{AN}{AC} = \dfrac{9}{12} = \dfrac34$: अनुपात बराबर, बिंदु एक ही
क्रम में, इसलिए थेल्स की विलोम से $(MN) \parallel (BC)$।
\end{solution}

\begin{solution}{exo:g9:thales:5}
$\dfrac{AM}{AB} = \dfrac46 = \dfrac23$ और
$\dfrac{AN}{AC} = \dfrac58$ हैं। मिलाकर जाँचो:
$\dfrac23 = \dfrac{16}{24}$ पर $\dfrac58 = \dfrac{15}{24}$ --- अनुपात
अलग हैं। अगर रेखाएँ \omterm{def:g6:lines:perp}{समांतर} होतीं, तो थेल्स प्रमेय उन्हें बराबर कर
देती; इसलिए $(MN)$ और $(BC)$ \omterm{def:g6:lines:perp}{समांतर} \emph{नहीं} हैं।
\end{solution}

\begin{solution}{exo:g9:thales:6}
मान लो $S$ छाया का सिरा है, $P$ व्यक्ति के सिर का ऊपरी बिंदु और $L$
खंभे का ऊपरी सिरा। व्यक्ति ($S$ से $3$ m दूर, $1.8$ m ऊँचा) और खंभा
($S$ से $3 + 2 = 5$ m दूर, ऊँचाई $h$) दो \omterm{def:g6:lines:perp}{समांतर} खड़े \omterm{def:g6:lines:objects}{रेखाखंड} हैं
जिन्हें रोशनी की किरण $(SL)$ काटती है: बिंदु $S$ से थेल्स लगाने पर
\[
\frac{1.8}{h} = \frac{3}{5},
\qquad\text{यानी}\quad
h = \frac{1.8 \times 5}{3} = 3 \text{ m।}
\]
\end{solution}

\begin{solution}{exo:g9:thales:7}
\emph{1.} असली दूरी: $6.8 \times 25\,000 = 170\,000$ cm $=
1.7$ km।

\emph{2.} \omterm{def:g6:measure:area}{क्षेत्रफल} $k^2$ के हिसाब से बदलते हैं: असली \omterm{def:g6:measure:area}{क्षेत्रफल}
$= 8 \times 25\,000^2$ cm$^2$ $= 8 \times 6.25 \times 10^8$ cm$^2$
$= 5 \times 10^9$ cm$^2$। और $1$ km$^2 = 10^{10}$ cm$^2$ होता है,
इसलिए यह $0.5$ km$^2$ है।
\end{solution}

\begin{solution}{exo:g9:thales:8}
\emph{1.} अनुपात $\dfrac{AM}{AB} = \dfrac{8}{12} = \dfrac23$ है।
थेल्स: $AN = \frac23 \times 15 = 10$ और $MN = \frac23 \times 18 = 12$।

\emph{2.} \omterm{def:g9:thales:scaling}{मापक गुणक} $k = \frac23$ है। \omterm{def:g6:measure:perimeter}{परिमाप}: $ABC$ का
$12 + 15 + 18 = 45$, और $AMN$ का
$8 + 10 + 12 = 30 = \frac23 \times 45$ --- \omterm{def:g6:measure:perimeter}{परिमाप} भी हर लंबाई की तरह
$k$ के हिसाब से बदलता है।
\end{solution}

\begin{solution}{exo:g9:thales:9}
$O$ के चारों ओर रेखाएँ $(AC)$ और $(BD)$ कटती हैं, और
$(AB) \parallel (CD)$ है: त्रिभुजों $OAB$ तथा $OCD$ वाला तितली
विन्यास। थेल्स देती है
\[
\frac{OA}{OC} = \frac{OB}{OD} = \frac{AB}{CD} = \frac46 = \frac23 .
\]
यानी $\dfrac{OA}{OC} = \dfrac23$, और बराबरी
$\dfrac{OA}{OC} = \dfrac{OB}{OD}$ ठीक यही कहती है कि $O$ दोनों
विकर्णों को एक ही अनुपात में काटता है।
\end{solution}

\begin{solution}{pb:g9:thales:1}
\textbf{1.} रचना से $[AB]$ पर दो बिंदु मिलते हैं; उन्हें $M_1$
($P_1$ से) और $M_2$ ($P_2$ से) कहो।

\textbf{2.} त्रिभुज $A P_3 B$ में, $P_1$ और $P_2$ से होकर
$(P_3 B)$ के \omterm{def:g6:lines:perp}{समांतर} खींची रेखाएँ भुजाओं को \omterm{def:g9:linfunc:linear}{समानुपाती} काटती हैं
(\cref{thm:g9:thales:direct}):
\[
\frac{A M_1}{AB} = \frac{A P_1}{A P_3} = \frac13,
\qquad
\frac{A M_2}{AB} = \frac{A P_2}{A P_3} = \frac23 .
\]
परकार ने $A P_1$, $P_1 P_2$, $P_2 P_3$ को बराबर बनाया था, इसलिए
अनुपात ठीक तिहाई हैं --- किरण का कोण और परकार का फैलाव चाहे जो भी
चुना गया हो।

\textbf{3.} पाँचवें हिस्सों के लिए: पाँच बराबर लंबाइयाँ नापो,
पाँचवें बिंदु को $B$ से जोड़ो, और चार \omterm{def:g6:lines:perp}{समांतर रेखाएँ} खींचो। किसी
\omterm{def:g6:lines:objects}{रेखाखंड} का $\frac35$ चाहिए तो वही चित्र, और $P_3$ से गुज़रती \omterm{def:g6:lines:perp}{समांतर}
रेखा जिस बिंदु पर काटती है वही $A$ से $[AB]$ के $\frac35$ पर है।

\textbf{4.} $\frac{AM}{MB} = \frac23$ का मतलब है
$AM = \frac25 AB$: किरण पर पाँच बराबर क़दम, $P_5$ को $B$ से जोड़ो,
और $P_2$ से गुज़रती \omterm{def:g6:lines:perp}{समांतर} रेखा $M$ पर निशान लगा देती है।

\textbf{5.} $\sqrt2 = 1.41421\ldots$ cm नापने में हमेशा गिने-चुने
दशमलव ही काम आते हैं: नापकर निकाला हुआ कोई भी ``तिहाई'' एक सन्निकट
मान ही रहेगा। थेल्स की रचना कोई संख्या पढ़ती ही नहीं: वह ठीक
$\frac{\sqrt2}{3}$ वाला बिंदु दे देती है --- ऐसे \omterm{def:g6:lines:objects}{रेखाखंड} के बराबर
हिस्से जिसे कोई पट्टी बता भी नहीं सकती
(\cref{pb:g9:sqrt:1})।

\textbf{6.} वस्तु का ऊपरी सिरा, छेद और तस्वीर का निचला सिरा एक ही
रेखा पर हैं, और वैसे ही वस्तु का निचला सिरा तथा तस्वीर का ऊपरी सिरा:
छेद पर कटती दो रेखाएँ, और वस्तु तथा तस्वीर वाली दीवार आपस में
\omterm{def:g6:lines:perp}{समांतर}। तितली में थेल्स:
\[
\frac{h}{H} = \frac{d}{D},
\qquad\text{यानी}\qquad
h = H \times \frac dD .
\]

\textbf{7.} $h = 12 \times \frac{0.20}{30} = 0.08$ m $= 8$ cm ---
और उलटी (किरणें छेद पर एक-दूसरे को काट जाती हैं)।

\textbf{8.} $H = h \times \frac Dd = 9 \times 10^{-3} \times
\frac{1.5 \times 10^{11}}{1} = 1.35 \times 10^{9}$ m। असली मान
$1.39 \times 10^9$ m है: जूते का एक डिब्बा सूरज को कुछ ही प्रतिशत की
ग़लती के भीतर नाप देता है।

\textbf{9.} चाँद: $\frac{3.5 \times 10^6}{3.8 \times 10^8}
\approx 0.0092$। सूरज: $\frac{1.39 \times 10^9}{1.5 \times
10^{11}} \approx 0.0093$। दोनों अनुपात --- यानी आसमान में दिखने वाले
आकार --- लगभग एक जैसे हैं: चाँद सूरज को \emph{बिलकुल} ढक सकता है,
\omterm{def:g5:solids:def}{किनारे} से \omterm{def:g5:solids:def}{किनारे} तक। यही ब्रह्मांडीय संयोग पूर्ण सूर्यग्रहण है।

\textbf{10.} हर किरण को उसी एक छेद से गुज़रना पड़ता है, इसलिए वस्तु
के ऊपरी सिरे से चली किरणें आगे \emph{नीचे} जाती हैं और निचले सिरे से
चलीं \emph{ऊपर}: तस्वीर उलट जाती है। छेद बड़ा करने पर तस्वीर की
थोड़ी-थोड़ी खिसकी हुई कई नक़लें एक साथ भीतर आती हैं और आपस में ऊपर-नीचे
पड़ जाती हैं: ज़्यादा चमकीली, पर फैली हुई।

\textbf{11.} $\frac{PQ}{MN} = \frac{OP}{OM} = \frac{7.5}{3}
= 2.5$, इसलिए $PQ = 10$; और $\frac{ON}{OQ} = \frac{OM}{OP} =
\frac{1}{2.5}$, इसलिए $ON = \frac{6}{2.5} = 2.4$।

\textbf{12.} अगर \omterm{def:g6:lines:perp}{समांतर} रेखा किसी एक भुजा के मध्यबिंदु से गुज़रे, तो
थेल्स का अनुपात $\frac12$ होता है, इसलिए वह दूसरी भुजा को भी
\emph{उसी के} मध्यबिंदु पर काटती है, और \omterm{def:g6:lines:perp}{समांतर} \omterm{def:g6:lines:objects}{रेखाखंड} आधार का \omterm{def:g3:division:half}{आधा}
नापता है: ठीक \cref{thm:g8:midpoints:direct}। थेल्स वही मध्यबिंदु
प्रमेय है, बस $\frac12$ अनुपात की बंदिश से आज़ाद।

\textbf{13.} त्रिभुज $ACD$ में रेखा $(EO)$ आधार $(CD)$ के \omterm{def:g6:lines:perp}{समांतर} है,
और $\frac{AO}{AC} = \frac{OA}{OA + OC} = \frac{4}{4 + 6} = \frac25$।
थेल्स: $\frac{EO}{DC} = \frac{AO}{AC} = \frac25$, इसलिए
$EO = \frac25 \times 6 = 2.4$।

\textbf{14.} त्रिभुज $BDC$ में उसी तरह $\frac{BO}{BD} = \frac25$ और
$OF = \frac25 \times 6 = 2.4$। इसलिए
\[
EF = EO + OF = 4.8
= \frac{2 \times 4 \times 6}{4 + 6} :
\]
यानी आधारों का \omterm{pb:g8:speed:1}{हरात्मक माध्य} --- आने-जाने के सफ़र वाला माध्य
(\cref{pb:g8:speed:1}), जो यहाँ एक समलंब में खिंचा हुआ है।

\textbf{15.} \omterm{def:g6:lines:perp}{समांतर}: $5$; गुणोत्तर: $\sqrt{24} \approx 4.90$;
हरात्मक: $4.8$। क्रम: $4.8 < 4.90\ldots < 5$, यानी हरात्मक $<$
गुणोत्तर $<$ \omterm{def:g6:lines:perp}{समांतर} --- और बराबरी सिर्फ़ तभी जब दोनों आधार बराबर हों।
चित्र में: मध्य रेखा $5$ नापती है, विकर्णों के कटान से गुज़रती \omterm{def:g6:lines:perp}{समांतर}
रेखा $4.8$, और समरूपता वाली काट $\sqrt{24}$ --- तीनों दोनों आधारों के
बीच इसी क्रम में एक के ऊपर एक बैठी हैं। असमिकाएँ
\cref{pb:g8:cosine:1} में सिद्ध हुई थीं (गुणोत्तर $\leq$ \omterm{def:g6:lines:perp}{समांतर},
वृत्त से भी और सर्वसमिका से भी) और \cref{pb:g8:speed:1} में (हरात्मक
$\leq$ \omterm{def:g6:lines:perp}{समांतर}); हरात्मक $\leq$ गुणोत्तर इन दोनों को जोड़ने से निकलता
है ($H \times A = G^2$, जो वहीं सिद्ध हुआ है)।

\end{solution}
